\chapter{背景}
\label{ch:background}

本章では，Arduino オーケストラというテーマを提案するに至った社会的・技術的な背景を整理する。

\section{社会的背景}
\label{sec:social_background}

\subsection{合奏・アンサンブルが持つ価値と参加障壁}

合奏は複数人が同じ時間軸の上で音を重ねることで初めて成立する表現形式であり，リズム感・他者への傾聴・即時的な調整といった音楽以外の能力も育む場として重視されてきた\cite{mext_music_curriculum}。一方で，オーケストラや吹奏楽のような生演奏のアンサンブルに参加するためには以下のような障壁が存在する。

\begin{itemize}
  \item \textbf{技能習得コスト}: 楽器演奏には長期間の訓練が必要であり，初学者が合奏に加わる機会は限られる。
  \item \textbf{同期の難しさ}: テンポや強弱を合わせる「息合わせ」は経験に依存し，初心者だけのアンサンブルでは成立しにくい。
  \item \textbf{物理的・経済的制約}: 楽器・防音・練習スペースなどのコストが高く，家庭や学校で気軽に取り組めない場合がある。
\end{itemize}

\subsection{技術を介した音楽体験の拡張ニーズ}

近年は電子楽器・DAW（Digital Audio Workstation）・リズムゲームなど，技能を前提としない音楽体験が広く普及している。特に新型コロナ禍以降，遠隔合奏・リモート演奏への関心が高まり，低遅延ネットワーク演奏や，スマートフォンを用いた参加型音楽コンテンツの研究・商用化が進んだ\cite{jacktrip}。こうした流れは「誰もが演奏に参加できる体験」を志向するものであり，合奏の敷居を下げる技術的取り組みが求められている。

\subsection{STEAM 教育との接続}

文部科学省が推進する STEAM 教育では，科学・技術と芸術を横断する題材が重視されている\cite{mext_steam}。マイコンとセンサを用いて「指揮」や「演奏」を電子的に再現する本テーマは，プログラミング・組込みシステム・音響処理・人間の動作解析といった要素を統合するため，工学系学生にとって学習テーマとしての価値が高く，将来的には教育用コンテンツへの展開も期待できる。

\section{技術的背景}
\label{sec:technical_background}

\subsection{低価格・無線対応マイコンの普及}

Arduino UNO R4 WiFi のような安価で Wi-Fi を搭載したマイコンが一般化したことで，複数ノードによる分散システムを学生レベルの予算でも構築できるようになった\cite{arduino_uno_r4}。特に UNO R4 WiFi は以下の特徴を持ち，本テーマと親和性が高い。

\begin{itemize}
  \item Arm Cortex-M4 コア（Renesas RA4M1）により従来の AVR よりも高速な信号処理が可能
  \item \textbf{IMU（LSM6DSOX）を基板上に内蔵}しており，外付けセンサなしで加速度・角速度を取得できる
  \item Wi-Fi モジュール（ESP32-S3）を搭載し，UDP/TCP による低遅延無線通信が容易
\end{itemize}

\subsection{IMU を用いたジェスチャ認識}

IMU（Inertial Measurement Unit）は加速度・角速度をミリ秒単位で取得できるため，スマートフォンやゲームコントローラにおける動作認識で広く利用されている\cite{lsm6dsox}。楽器演奏や指揮動作のような周期性のあるモーションは，閾値ベースのピーク検出や機械学習モデルを用いて比較的高い精度で認識できることが報告されている\cite{nintendo_wiimusic,imu_gesture_survey}。指揮棒の振り下ろし・振り幅・振動といった特徴量は加速度ベクトルのノルムや主軸成分から抽出可能である。

\subsection{低遅延無線通信プロトコル}

音楽演奏では人間が許容できる同期ズレは数十ミリ秒程度とされる\cite{music_latency}。Wi-Fi 上で動作する UDP は TCP よりも再送制御のオーバーヘッドが小さく，低遅延・リアルタイム用途に向くため，本テーマのような複数楽器への一斉配信に適する\cite{udp_realtime}。同期精度は，タイムスタンプ付きパケット・NTP 的な時刻同期・バッファリングなどの組み合わせで確保される。

\subsection{クリエイティブコーディング環境による音響合成}

Processing や openFrameworks のような環境は，GUI やサウンド合成を比較的短いコードで記述でき，ビジュアルとオーディオを同時に扱える\cite{processing_sound}。マイコン単体では PWM による単音程度しか鳴らせないが，マイコン側を「楽譜データ送出器」として扱い，音色合成は PC 側の Processing に任せる構成にすることで，金管楽器の倍音構造や ADSR エンベロープのような表現力の高い音色合成が実現できる。

\section{本テーマの位置付け}
\label{sec:positioning}

以上より，本テーマ「Arduino オーケストラ」は次の 3 つの流れの接点に位置する。

\begin{enumerate}
  \item 合奏への参加障壁を下げる \textbf{社会的要請}
  \item 無線マイコン・IMU・低遅延通信の \textbf{技術的成熟}
  \item STEAM 教育・エンタテインメントとしての \textbf{応用可能性}
\end{enumerate}

次章では，これらに関連する既存システムを整理し，本提案の差別化点を明確化する。
